\documentclass[../DoAn.tex]{subfiles}
\begin{document}
\label{ch:bai_toan}
% =============================================================
% CHƯƠNG 3 - BÀI TOÁN NHÚNG MẠNG ẢO ĐA MIỀN
% =============================================================

Chương này hình thức hoá bài toán nhúng mạng ảo đa miền mà đồ án hướng tới, rồi điểm qua các nghiên cứu liên quan. Phần đầu mô tả mô hình mạng vật lý đa miền và mạng ảo, hai pha ánh xạ nút và liên kết, hàm mục tiêu cùng các độ đo, phát biểu bài toán ở dạng vận hành (đầu vào, quyết định, mục tiêu trong kịch bản trực tuyến) và mô hình hoá đầy đủ dưới dạng quy hoạch nguyên tuyến tính (ILP). Phần cuối khảo sát các hướng tiếp cận hiện có và định vị đóng góp của đồ án.

\section{Mạng vật lý đa miền và mạng ảo}
\label{sec:vne_model}

\textbf{Mạng vật lý} (substrate network) được biểu diễn bằng đồ thị vô hướng có trọng số
\[
    G^s = (N^s, E^s),
\]
trong đó $N^s$ là tập các nút vật lý và $E^s$ là tập các liên kết. Mỗi nút $n^s \in N^s$ được đặc trưng bởi dung lượng CPU $C(n^s)$, đơn giá CPU $P(n^s)$ và độ trễ xử lý $D(n^s)$. Mỗi liên kết $e^s = (u^s, v^s) \in E^s$ được đặc trưng bởi băng thông $B(e^s)$, đơn giá băng thông $P(e^s)$ và độ trễ truyền dẫn $D(e^s)$. Tại thời điểm bất kỳ, tài nguyên còn lại của mỗi nút và liên kết là $C^{\text{avail}}(n^s) \leq C(n^s)$ và $B^{\text{avail}}(e^s) \leq B(e^s)$.

\textbf{Cấu trúc đa miền.} Khác với bài toán VNE đơn miền, mạng vật lý ở đây được chia thành một tập miền quản trị $\mathcal{D}=\{d_1,\ldots,d_D\}$. Mỗi miền $d$ là một đồ thị con với các \textit{liên kết nội miền}; các miền nối với nhau qua các \textit{liên kết liên miền} (inter-domain link) giữa những \textit{nút biên} (boundary node). Một giả định cốt lõi của cấu hình đa miền là thông tin chi tiết của các nút bên trong một miền \textit{không hiển thị toàn cục}: chỉ thông tin tổng hợp ở mức miền và các liên kết liên miền được trao đổi giữa các miền. Đây là điểm khác biệt chính so với VNE đơn miền và là yếu tố mà phương pháp đề xuất phải tôn trọng khi thiết kế bộ mã hoá (Chương~\ref{ch:de_xuat}).

Một \textbf{yêu cầu mạng ảo} (Virtual Network Request - VNR) bao gồm đồ thị mạng ảo và các thông tin thời gian:
\[
    \text{VNR} = (G^v, t_a, t_l), \quad G^v = (N^v, E^v),
\]
với $t_a$ là thời điểm đến và $t_l$ là thời gian sống. Mỗi nút ảo $n^v \in N^v$ có yêu cầu CPU $c(n^v)$, và mỗi liên kết ảo $e^v = (u^v, v^v) \in E^v$ có yêu cầu băng thông $b(e^v)$. Trong cấu hình đa miền, mỗi nút ảo còn có thể bị ràng buộc bởi tập \textit{miền cho phép} $\mathcal{D}(n^v) \subseteq \mathcal{D}$ (ví dụ do yêu cầu vị trí địa lý hay pháp lý): nút ảo chỉ được ánh xạ vào nút vật lý thuộc một trong các miền này.

\section{Hai pha: ánh xạ nút và ánh xạ liên kết}
\label{sec:two-phase}

Một lời giải nhúng cho VNR là một cặp ánh xạ $(f_N, f_E)$ trong đó:
\begin{itemize}
\item \textbf{Ánh xạ nút} $f_N : N^v \to N^s$ thoả mãn (i) ràng buộc tài nguyên $c(n^v) \leq C^{\text{avail}}(f_N(n^v))$ với mọi $n^v$, (ii) ràng buộc đơn ánh: $f_N(u^v) \neq f_N(v^v)$ với $u^v \neq v^v$, và (iii) ràng buộc miền $f_N(n^v) \in \mathcal{D}(n^v)$ nếu có.
\item \textbf{Ánh xạ liên kết} $f_E : E^v \to \mathcal{P}(E^s)$ ánh xạ mỗi liên kết ảo tới một tập đường đi vật lý nối $f_N(u^v)$ và $f_N(v^v)$, với tổng băng thông trên các đường đi bằng $b(e^v)$ và không vượt quá băng thông còn lại của bất kỳ liên kết vật lý nào.
\end{itemize}

Do tính chất NP-Hard của bài toán, các phương pháp thực tế thường giải hai pha một cách phối hợp nhưng không đồng thời tối ưu: pha ánh xạ nút sử dụng một chính sách xếp hạng/lựa chọn, sau đó pha ánh xạ liên kết dùng đường đi ngắn nhất hoặc multi-path để hoàn thiện. Đồ án này đi theo hướng đó, với cải tiến là sử dụng học tăng cường để học chính sách chọn ứng viên thay cho heuristic cứng, đồng thời lưu ý chi phí định tuyến xuyên miền.

\section{Hàm mục tiêu và độ đo hiệu năng}
\label{sec:objective}

Đối với mỗi VNR được nhúng thành công, \textbf{doanh thu} (revenue) và \textbf{chi phí} (cost) được định nghĩa:
\begin{equation}
    \mathcal{R}(\text{VNR}) = \sum_{n^v \in N^v} c(n^v) + \sum_{e^v \in E^v} b(e^v),
    \label{eq:revenue}
\end{equation}
\begin{equation}
    \mathcal{C}(\text{VNR}) = \sum_{n^v \in N^v} c(n^v) + \sum_{e^v \in E^v} b(e^v) \cdot \ell(f_E(e^v)),
    \label{eq:cost}
\end{equation}
trong đó $\ell(\cdot)$ là tổng số chặng vật lý của tập đường đi ánh xạ liên kết ảo. Hai độ đo hiệu năng phổ biến nhất là:
\begin{itemize}
\item \textbf{Tỉ lệ chấp nhận} (Acceptance Ratio - AR): tỉ số giữa số VNR được nhúng thành công và tổng số VNR đến trong khoảng thời gian khảo sát.
\item \textbf{Tỉ số doanh thu trên chi phí} (Revenue-to-Cost - R/C): được định nghĩa cho mỗi VNR là $\mathcal{R}(\text{VNR}) / \mathcal{C}(\text{VNR})$, đo mức độ ``tiết kiệm tài nguyên'' khi nhúng. Ở mức hệ thống, R/C tổng được tính bằng tỉ số của tổng doanh thu và tổng chi phí trên toàn bộ VNR đã nhúng.
\end{itemize}

Ngoài AR và R/C, các độ đo bổ sung gồm thời gian chạy trung bình mỗi VNR, độ trễ trung bình của các đường đi, và độ ổn định của reward trong quá trình huấn luyện. Đồ án này sử dụng đồng thời cả bốn độ đo để so sánh.

\section{Phát biểu bài toán}
\label{sec:problem-statement}

Dựa trên các thành phần đã định nghĩa, bài toán đồ án giải quyết được phát biểu như sau.

\textbf{Đầu vào.} Một mạng vật lý đa miền $G^s = (N^s, E^s)$ với tập miền quản trị $\mathcal{D}$, trong đó thông tin nội miền chỉ hiển thị ở mức tổng hợp; và một \textit{dòng} yêu cầu mạng ảo $\{\text{VNR}_1, \text{VNR}_2, \ldots\}$ đến tuần tự theo thời gian (kịch bản trực tuyến), mỗi $\text{VNR}_i = (G^v_i, t_a^i, t_l^i)$ kèm yêu cầu CPU, băng thông và ràng buộc miền cho phép $\mathcal{D}(n^v)$.

\textbf{Quyết định.} Tại thời điểm mỗi VNR đến, dựa trên tài nguyên còn lại của hạ tầng, phải quyết định \textit{chấp nhận} hay \textit{từ chối} VNR đó; nếu chấp nhận, phải xác định một lời giải nhúng $(f_N, f_E)$ thoả mãn đồng thời ràng buộc tài nguyên, đơn ánh, băng thông và ràng buộc miền. Quyết định mang tính trực tuyến và không thể thu hồi: lời giải cho $\text{VNR}_i$ được đưa ra mà không biết trước các VNR tương lai, và tài nguyên đã cấp phát bị giữ đến khi VNR hết thời gian sống $t_l^i$ rồi mới được giải phóng.

\textbf{Mục tiêu.} Tối đa hoá hiệu năng tích luỹ trên toàn dòng yêu cầu, ưu tiên \textit{tỉ lệ chấp nhận} (AR) và \textit{tỉ số doanh thu trên chi phí} (R/C), đồng thời tôn trọng đặc thù đa miền: hạn chế thông tin toàn cục và chi phí định tuyến xuyên miền phải được tính ngay tại pha ra quyết định.

Đây là một bài toán tối ưu tổ hợp NP-Hard (Mục~\ref{sec:ilp}), lại đặt trong môi trường trực tuyến với thông tin hạn chế, nên các lời giải thực tế hướng tới chính sách ra quyết định nhanh và tổng quát hoá tốt thay vì lời giải tối ưu tuyệt đối cho từng VNR.

\section{Mô hình hoá bài toán dưới dạng ILP}
\label{sec:ilp}

Để có mô hình hoá đầy đủ, đặt biến nhị phân $x_{n^v, n^s} \in \{0,1\}$ bằng $1$ nếu nút ảo $n^v$ được ánh xạ tới nút vật lý $n^s$, và biến $y^{u^v v^v}_{u^s v^s} \in [0, 1]$ biểu diễn tỉ lệ băng thông của liên kết ảo $(u^v, v^v)$ đi qua liên kết vật lý $(u^s, v^s)$. Với mục tiêu cực đại hoá tỉ số R/C, bài toán có thể viết:
\begin{equation}
    \max \quad \frac{\mathcal{R}(\text{VNR})}{\mathcal{C}(\text{VNR})}
\end{equation}
với các ràng buộc:
\begin{align}
    \sum_{n^s \in N^s} x_{n^v, n^s} &= 1, \quad \forall n^v \in N^v, \label{eq:c1}\\
    \sum_{n^v \in N^v} x_{n^v, n^s} \cdot c(n^v) &\leq C^{\text{avail}}(n^s), \quad \forall n^s \in N^s, \label{eq:c2}\\
    \sum_{n^v \in N^v} x_{n^v, n^s} &\leq 1, \quad \forall n^s \in N^s, \label{eq:c3}\\
    \sum_{(u^s, v^s) \in E^s} y^{u^v v^v}_{u^s v^s} \cdot b(u^v, v^v) &\leq B^{\text{avail}}(u^s, v^s), \quad \forall (u^s, v^s), \label{eq:c4}\\
    \sum_{v^s : (u^s, v^s) \in E^s} (y^{u^v v^v}_{u^s v^s} - y^{u^v v^v}_{v^s u^s}) &= x_{u^v, u^s} - x_{v^v, u^s}, \quad \forall u^s, (u^v, v^v). \label{eq:c5}
\end{align}
Ràng buộc \eqref{eq:c1} đảm bảo mỗi nút ảo được ánh xạ tới đúng một nút vật lý; \eqref{eq:c2} đảm bảo không vượt CPU; \eqref{eq:c3} đảm bảo tính đơn ánh; \eqref{eq:c4} đảm bảo không vượt băng thông; và \eqref{eq:c5} là ràng buộc bảo toàn luồng (multi-commodity flow) cho ánh xạ liên kết. Trong cấu hình đa miền có thể bổ sung ràng buộc $x_{n^v,n^s}=0$ với mọi $n^s$ không thuộc miền cho phép $\mathcal{D}(n^v)$. Dạng ILP này có $|N^v| \cdot |N^s|$ biến nhị phân nút và $|E^v| \cdot |E^s|$ biến liên kết; do không gian nghiệm tăng theo cấp số mũ, việc giải tối ưu trở nên không khả thi với mạng vật lý cỡ trăm nút trở lên (phân tích độ phức tạp chi tiết được trình bày ở Mục~\ref{sec:complexity-time} của Chương~\ref{ch:de_xuat}). Đây chính là lý do đồ án theo đuổi hướng xấp xỉ bằng học tăng cường thay vì giải ILP trực tiếp.

\section{Nghiên cứu liên quan}
\label{sec:approaches}

Bài toán nhúng mạng ảo đã được nghiên cứu rộng rãi trong hơn hai thập kỷ; các khảo sát tổng quan~\cite{fischer2013vne_survey, cao2019vne_survey} phân loại lời giải thành ba nhóm chính. Nhóm \textit{lời giải tối ưu và cận tối ưu} phát biểu bài toán dưới dạng quy hoạch nguyên (hỗn hợp) tuyến tính~\cite{andersen2002theoretical, zhu2006algorithms}; tiêu biểu là ViNE~\cite{chowdhury2009vine}, phối hợp pha ánh xạ nút và liên kết qua nới lỏng tuyến tính (LP-relaxation) rồi làm tròn, cho lời giải chất lượng cao và cận trên hữu ích nhưng độ phức tạp cấp số mũ khiến chúng chỉ khả thi trên tô-pô rất nhỏ và không phù hợp kịch bản trực tuyến quy mô lớn.

Nhóm \textit{heuristic} xếp hạng nút vật lý theo những độ đo tô-pô đơn giản rồi ánh xạ tham lam: NodeRank lấy cảm hứng từ PageRank~\cite{cheng2011vne_nr}, Global Resource Capacity (GRC) tổng hợp tài nguyên lan truyền trên đồ thị~\cite{gong2014grc}, ánh xạ phân tán~\cite{houidi2008distributed}, phát hiện đồng cấu đồ thị con~\cite{lischka2009virtual}, hay cho phép chia đường và di trú~\cite{yu2008rethinking}. Heuristic được so sánh trực tiếp trong đồ án là MP-VNE~\cite{mp_vne}, ước lượng chi phí từng cặp $(n^v, n^s)$ bằng công thức tổng hợp giữa chi phí CPU và băng thông trung bình của láng giềng (gọi là PreCost) rồi chọn các nút có chi phí ước lượng thấp nhất làm ứng viên, với pha ánh xạ liên kết dùng đường đi ngắn nhất có thể mở rộng đa đường. Heuristic đơn giản, nhanh và chạy được trên hạ tầng lớn, nhưng phụ thuộc nặng vào tri thức chuyên gia, khó thích nghi khi phân phối yêu cầu thay đổi và dễ rơi vào tối ưu cục bộ. Nhóm \textit{metaheuristic} như thuật toán di truyền, tối ưu hoá bầy hạt (PSO) và mô phỏng tôi luyện biểu diễn mỗi lời giải như một cá thể với hàm thích nghi là chi phí nhúng rồi tiến hoá qua các thế hệ; chẳng hạn MP-VNE~\cite{mp_vne} kết hợp xếp hạng PreCost với một vòng PSO. Nhóm này cho chất lượng khá hơn heuristic nhưng chi phí tính toán cao và thiếu đảm bảo lý thuyết về chất lượng lời giải.

Gần đây, học tăng cường (Reinforcement Learning - RL) coi VNE như một quá trình quyết định tuần tự và học chính sách ánh xạ end-to-end từ tín hiệu thưởng. Mijumbi và cộng sự~\cite{mijumbi2014rl} là một trong những công trình tiên phong áp dụng RL cho quản lý tài nguyên mạng ảo động; với sự phát triển của mạng nơ-ron đồ thị~\cite{kipf2017gcn, velickovic2018gat, wu2021gnn_survey}, nhiều phương pháp dùng GNN để mã hoá tô-pô mạng vật lý như A3C-GCN~\cite{yao2018a3c_gcn}, DeepViNE~\cite{dolati2019deepvine}, NeuroViNE~\cite{blenk2018neurovine} và FlagVNE~\cite{flagvne2024}; một số công trình khác dùng tìm kiếm cây Monte-Carlo~\cite{haeri2018mcts} hoặc cơ chế chú ý cho bài toán định tuyến tổ hợp~\cite{kool2019attention_co}. Mặc dù nhiều tiềm năng, các phương pháp RL hiện hành cho VNE còn ba hạn chế chính, cũng là động lực trực tiếp cho phương pháp đề xuất: (i) phụ thuộc mạnh vào \textit{reward shaping} chính xác, dễ bất ổn định khi reward thưa; (ii) chưa khai thác tốt tri thức tô-pô có sẵn từ các heuristic mạnh, khiến tác nhân tốn nhiều mẫu để khám phá lại các quy luật đơn giản; và (iii) không gian hành động lớn (chọn một trong hàng trăm nút vật lý cho mỗi nút ảo) làm chính sách khó hội tụ ở tô-pô lớn.

Phần lớn các công trình trên giả định hạ tầng \textit{đơn miền} với thông tin tô-pô và tài nguyên hiển thị toàn cục. Trong thực tế, hạ tầng thường được chia thành nhiều miền quản trị thuộc các nhà cung cấp khác nhau, nơi thông tin nội miền không công khai toàn cục mà chỉ trao đổi ở mức tổng hợp~\cite{cao2019vne_survey}, đặt thêm hai thách thức: tôn trọng ràng buộc miền của từng nút ảo và tính tới chi phí định tuyến xuyên miền ngay khi ra quyết định ánh xạ. Đồ án tập trung vào cấu hình đa miền này, kế thừa heuristic PreCost của MP-VNE làm baseline nhưng thay chính sách xếp hạng/chọn ứng viên thủ công bằng một chính sách học được tôn trọng cấu trúc đa miền (Chương~\ref{ch:de_xuat}). Bảng~\ref{tab:approach-compare} tổng hợp ưu nhược của các nhóm tiếp cận.

\begin{table}[h]
\centering
\caption{So sánh các hướng tiếp cận VNE}
\label{tab:approach-compare}
\resizebox{\textwidth}{!}{%
\begin{tabular}{l l l l}
\hline
\textbf{Tiêu chí} & \textbf{Heuristic} & \textbf{Metaheuristic} & \textbf{Học tăng cường} \\
\hline
Chất lượng lời giải & Trung bình & Khá & Tốt (khi huấn luyện ổn định) \\
Thời gian suy diễn & Rất nhanh & Chậm & Nhanh \\
Khả năng tổng quát hoá & Thấp & Trung bình & Cao \\
Phụ thuộc tri thức chuyên gia & Cao & Trung bình & Thấp \\
Chi phí huấn luyện & Không có & Không có & Cao \\
\hline
\end{tabular}%
}
\end{table}

% Kết chương
Chương này đã phát biểu hình thức bài toán nhúng mạng ảo đa miền (mô hình mạng vật lý đa miền và mạng ảo, hai pha ánh xạ, hàm mục tiêu, phát biểu bài toán trực tuyến và mô hình hoá ILP) và khảo sát các nghiên cứu liên quan. Chương tiếp theo trình bày thuật toán đề xuất CARL-VNE.

\end{document}
